\documentclass[11pt, a4paper, sans]{moderncv}
\moderncvstyle{classic}
\moderncvcolor{blue}

\usepackage[scale=0.90, top=1.5cm, bottom=1.5cm]{geometry}
\setlength{\hintscolumnwidth}{3cm}

\name{Leonardo}{Santana}
\title{Desenvolvedor Fullstack Pleno}
\email{leosantanabr@gmail.com}
\phone[mobile]{+55~75~98146~0783}
\address{Rua Volta Do Tanque, 165}{44400-000 Nazaré -- BA}
\social[linkedin]{banzak}
\social[github]{banzak1}

\begin{document}
\makecvtitle
\vspace{-8mm}

\section{Resumo Profissional}
\cvitem{}{Desenvolvedor Fullstack Pleno com 4+ anos de experiência em sistemas de missão crítica --- do back-end Node.js/Java ao front-end Angular/React, passando por mensageria com Apache Kafka, WebSocket e bancos SQL/NoSQL. Construí sistemas de processamento em tempo real que atualizam dados a cada 1 segundo e APIs de alta frequência com Socket.IO. No dia a dia, uso Claude Code e Spec-Driven Development para acelerar o ciclo de entrega com qualidade.}

\section{Competências}
\cvitem{\textbf{Linguagens}}{Java, JavaScript, TypeScript, Node.js, Go}
\cvitem{\textbf{Front-end}}{Vue.js, Angular, React}
\cvitem{\textbf{Back-end}}{Spring Boot, Node.js, Express.js, APIs RESTful, Microsserviços, Clean Code, SOLID}
\cvitem{\textbf{Infra \& Dados}}{PostgreSQL, Oracle, MongoDB, ElasticSearch, Apache Kafka, Docker, Socket.IO}
\cvitem{\textbf{Ferramentas \& Processos}}{Git, JUnit, Jenkins, Jira, Scrum, Claude Code, Spec-Driven Development}

\section{Experiência}
\cventry{Out 2025--Atual}{Desenvolvedor Fullstack Pleno}{DB1 Global Software}{Maringá}{}{
\begin{itemize}
    \item Responsável pelo front-end do Omni Conectado, plataforma PWA (Angular, SCSS, Jest) para simulação de financiamento de veículos usados --- leves, motos e, em breve, pesados --- utilizada por lojistas de todo o Brasil.
    \item Atuação direta em todo o ciclo de entrega: criação de cards no Jira, desenvolvimento, code review rigoroso com Git Flow e deploy via Jenkins.
    \item Defesa ativa das implantações em reuniões de GMUD, garantindo alinhamento entre desenvolvimento e operações.
    \item Suporte e manutenção do sistema legado de financiamento de veículos pesados e pessoa física, garantindo continuidade do negócio enquanto a nova plataforma evolui.
\end{itemize}}

\cventry{Mai 2022--Jun 2025}{Analista de Desenvolvimento JR}{Solinftec}{Araçatuba}{}{
\begin{itemize}
    \item Implementei mapa de nuvens em Node.js com atualização a cada 10 segundos, renderizando variações de intensidade de chuva (mm) por cor --- imagens processadas via AWS e exibidas com zoom progressivo utilizando Socket.IO para comunicação em tempo real.
    \item Construí APIs em Go para dados climáticos (umidade, precipitação, temperatura, vento) consumidos via Kafka, com filtros por cliente, zona ou talhão e gráficos históricos utilizando PostGIS para geolocalização.
    \item Desenvolvi e mantive o sistema de visualização em tempo real de maquinários agrícolas (Java/Spring Boot + Kafka), processando dados de sensores (velocidade, consumo, temperatura, operador) a cada 1 segundo e persistindo em Oracle --- base para relatórios operacionais no PowerBI.
\end{itemize}}

\cventry{Nov 2021--Mai 2022}{Trainee Dev}{Solinftec}{Araçatuba}{}{
\begin{itemize}
    \item Imersão prática em Java, Vue.js, Node.js, Kafka, Oracle e PostgreSQL com projetos desafio como: sistema de mapa com movimentação de maquinários em tempo real e plataforma de compra/venda de ações fictícias com atualização a cada segundo via WebSocket + Kafka.
\end{itemize}}

\cventry{Jun 2017--Dez 2018}{Desenvolvedor Web}{Engenhe JR}{Cruz das Almas}{}{
\begin{itemize}
    \item Desenvolvimento de sites corporativos para clientes como TFR Agrícola e FederalCar.
\end{itemize}}

\section{Formação}
\cventry{Jan 2020--Ago 2022}{Análise e Desenvolvimento de Sistemas}{Estácio}{Nazaré}{}{}
\cventry{Set 2016--Dez 2019}{Bacharelado em Ciências Exatas e Tecnológicas}{Universidade Federal do Recôncavo da Bahia}{Cruz das Almas}{}{}

\section{Idiomas}
\cvitem{Português}{Nativo}
\cvitem{Inglês}{Técnico (leitura de documentação, escrita de código) --- conversação em desenvolvimento}

\end{document}
